\documentclass[12pt, a4paper]{article}
\usepackage[utf8]{inputenc}
\usepackage{geometry}
\geometry{margin=1in}
\usepackage{amsmath, amssymb, amsthm}
\usepackage{graphicx}
\usepackage{hyperref}
\usepackage{listings}
\usepackage{xcolor}

\definecolor{codegreen}{rgb}{0,0.6,0}
\definecolor{codegray}{rgb}{0.5,0.5,0.5}
\definecolor{codepurple}{rgb}{0.58,0,0.82}
\definecolor{backcolour}{rgb}{0.95,0.95,0.92}

\lstdefinestyle{mystyle}{
    backgroundcolor=\color{backcolour},   
    commentstyle=\color{codegreen},
    keywordstyle=\color{magenta},
    numberstyle=\tiny\color{codegray},
    stringstyle=\color{codepurple},
    basicstyle=\ttfamily\footnotesize,
    breakatwhitespace=false,         
    breaklines=true,                 
    captionpos=b,                    
    keepspaces=true,                 
    numbers=left,                    
    numbersep=5pt,                  
    showspaces=false,                
    showstringspaces=false,
    showtabs=false,                  
    tabsize=2
}

\lstset{style=mystyle}

\title{\textbf{Monte Carlo Option Pricing Engine}\\
\Large Bridging Theoretical Finance with Live Market Data}
\author{Midterm Project Report}
\date{\today}

\begin{document}

\maketitle

\begin{abstract}
This report details the theoretical foundation and practical implementation of a Monte Carlo Option Pricing Engine. We explore the genesis of the landmark Black-Scholes model, diving deeply into Geometric Brownian Motion (GBM) and Ito's Lemma. We also examine the critical assumption of constant volatility---an assumption that contributed to the historic collapse of Long-Term Capital Management (LTCM). Finally, we demonstrate the architectural implementation of the pricing engine, including how we extract live market volatility and risk-free rates using Python, setting the stage for advanced stochastic modeling in future iterations.
\end{abstract}

\tableofcontents
\newpage

\section{Introduction}
Financial markets are complex, continuous-time stochastic systems. In 1973, Fischer Black and Myron Scholes published their seminal paper, \textit{"The Pricing of Options and Corporate Liabilities,"} fundamentally altering the landscape of quantitative finance by formally connecting thermodynamics physics equations to financial derivatives.

The goal of this project is to construct a software engine capable of replicating these complex mathematical truths and applying them to live, real-time options data. This midterm report acts as the bridge between the academic theory (GBM, Ito's Lemma) and the software engineering required to build a functional Single Option Analysis tool.

\section{The Genesis: The Story of the Black-Scholes Model}
Prior to 1973, investors knew that an option's premium was related to the underlying stock, but they lacked a rigorous framework to determine the precise "Fair Value."

Fischer Black and Myron Scholes (alongside Robert Merton) recognized that if an investor continuously delta-hedged an option by buying and selling the underlying stock, they could construct a completely riskless portfolio. The fundamental economic principle of \textit{no-arbitrage} dictates that a portfolio with zero risk must earn exactly the risk-free rate of return.

This realization led to the derivation of the **Black-Scholes Partial Differential Equation (PDE)**:

\begin{equation}
\frac{\partial V}{\partial t} + \frac{1}{2}\sigma^2 S^2 \frac{\partial^2 V}{\partial S^2} + rS \frac{\partial V}{\partial S} - rV = 0
\end{equation}

\subsection{Deconstructing the PDE Variables}
\begin{itemize}
    \item \textbf{$V(S, t)$ (Value):} The theoretical fair price of the option.
    \item \textbf{$S$ (Spot Price):} The current price of the underlying asset. The derivatives concerning $S$ track how the option's value changes (first derivative = \textit{Delta}, second derivative = \textit{Gamma}).
    \item \textbf{$t$ (Time):} The ticking clock. The term $\frac{\partial V}{\partial t}$ represents \textit{Theta} decay.
    \item \textbf{$r$ (Risk-Free Rate):} The theoretical return of a zero-risk investment, traditionally anchored to US Treasury yields.
    \item \textbf{$\sigma$ (Volatility):} The crux of the model. The formulation explicitly assumes this parameter is \textit{constant} and measurable.
\end{itemize}

\section{Stochastic Calculus: Deep Dive into GBM}
To arrive at the Black-Scholes PDE, we must first model the behavior of the stock market itself. Stock prices do not move in straight lines; they follow a random walk with an upward drift, mathematically represented as **Geometric Brownian Motion (GBM)**.

The Stochastic Differential Equation (SDE) for a stock following GBM is:
\begin{equation}
dS = \mu S \, dt + \sigma S \, dW
\end{equation}
Where:
\begin{itemize}
    \item $dS$: the infinitesimal change in stock price.
    \item $\mu$: the expected return (drift).
    \item $\sigma$: the volatility.
    \item $dW$: standard Brownian motion (a Wiener process), introducing the random "noise".
\end{itemize}

\subsection{Ito's Lemma}
Because $dW$ introduces randomness with non-differentiable paths, standard calculus (Newtonian calculus) fails. We instead use **Ito's Lemma**, which acts as the stochastic equivalent of the Chain Rule. 

For a function $V(S, t)$, Ito's Lemma states:
\begin{equation}
dV = \left( \frac{\partial V}{\partial t} + \mu S \frac{\partial V}{\partial S} + \frac{1}{2} \sigma^2 S^2 \frac{\partial^2 V}{\partial S^2} \right) dt + \sigma S \frac{\partial V}{\partial S} \, dW
\end{equation}

By equating the return of the hedged portfolio to the risk-free rate, the random $dW$ term is algebraically eliminated, resulting in the deterministic Black-Scholes PDE shown in Equation 1.

\section{Monte Carlo Assumptions \& Implementation}
While the Black-Scholes formula is an elegant closed-form solution for European options, it struggles with complex, path-dependent exotic options. To create a robust and flexible engine, we utilized **Monte Carlo simulations**.

\subsection{Underlying Assumptions}
Our Monte Carlo engine operates under specific assumptions derived directly from the GBM framework:
\begin{enumerate}
    \item \textbf{Log-Normal Distribution:} Stock prices cannot drop below zero. Consequently, continuous compounding dictates that stock returns follow a normal distribution, meaning the stock prices themselves follow a log-normal distribution.
    \item \textbf{Constant Volatility:} The market "wiggles" at a constant, predictable rate ($\sigma$).
    \item \textbf{Frictionless Markets:} There are no transaction costs, Bid-Ask spreads, or restrictions on short selling.
\end{enumerate}

Using the exact solution to the GBM stochastic differential equation (SDE), terminal prices ($S_T$) are simulated via:
\begin{equation}
S_T = S_0 \exp\left( \left( r - \frac{\sigma^2}{2} \right) T + \sigma \sqrt{T} Z \right)
\end{equation}
where $Z \sim \mathcal{N}(0, 1)$ is a standard normal random variable. By simulating 10,000+ potential future paths and averaging the discounted payoffs, the Law of Large Numbers ensures our Monte Carlo engine converges flawlessly to the Black-Scholes price.

\section{Bridging Theory with Live Code}
Academic formulas often fail when not married to accurately sourced market parameters. A primary focus of the Single Option Analysis module was dynamically fetching real-world data to calculate $r$ and $\sigma$.

\subsection{Gathering the Risk-Free Rate ($r$)}
In academia, $r$ is often hardcoded to 5\%. In reality, the risk-free rate fluctuates daily based on central bank policy. Our system actively scrapes the yield of the 13-Week US Treasury Bill (\texttt{\textasciicircum IRX}) as the proxy for $r$:

\begin{lstlisting}[language=Python, caption=Fetching the Live Risk-Free Rate]
def get_risk_free_rate():
    try:
        irx = yf.Ticker("^IRX")
        hist = irx.history(period="5d")
        if not hist.empty:
            # ^IRX is returned as a percentage (e.g., 4.54 for 4.54%)
            latest_yield = float(hist['Close'].iloc[-1])
            return latest_yield / 100.0
    except Exception as e:
        print(f"Error fetching IRX: {e}")
    # Fallback default if API fails
    return 0.05 
\end{lstlisting}

\subsection{Calculating Historical Volatility ($\sigma$)}
Volatility ($\sigma$) cannot be observed directly; it must be derived. Our engine downloads 1 year of daily historical prices for the target underlying asset and computes the annualized standard deviation of its logarithmic returns:

\begin{lstlisting}[language=Python, caption=Calculating Annualized Historical Volatility, label=lst:vol]
def get_spot_and_vol(ticker_symbol):
    tk = yf.Ticker(ticker_symbol)
    history = tk.history(period="1y")

    # The current Spot Price (S0)
    spot = float(history['Close'].iloc[-1])
    
    # Calculate daily log returns
    log_returns = np.log(history['Close'] / history['Close'].shift(1)).dropna()
    
    # Annualize the volatility (assuming 252 trading days/year)
    hist_vol = float(log_returns.std() * np.sqrt(252))
    
    return {'spot': spot, 'historical_vol': hist_vol}
\end{lstlisting}

\section{The Fatal Flaw: The Story of LTCM}
A critical element of this project is understanding the limitations of the models we are building. The assumption of **constant volatility** and log-normal distributions means that massive market crashes (Fat-Tail events) are modeled as statistically virtually impossible.

This exact mathematical assumption led to the historic collapse of **Long-Term Capital Management (LTCM)**. Founded by a "Dream Team" of Wall Street traders and bolstered by board members Myron Scholes and Robert Merton (the architects of the Black-Scholes model who won the 1997 Nobel Prize), LTCM utilized these continuous-time models to hyper-leverage billions of dollars in arbitrage trades.

In August 1998, the Russian government unexpectedly defaulted on its domestic sovereign debt (the GKO crisis). Global markets panicked. The markets did not behave as smooth, continuous Brownian Motion; instead, they "gapped" violently downwards, creating a discontinuous shock to the financial system. Because their models ignored the possibility of discontinuous jumps and fluctuating volatility, LTCM lost \$4.6 billion in a matter of months, requiring an unprecedented Federal Reserve bailout to prevent a global financial meltdown.

\section{Conclusion and Future Scope}
The initial phase of this project successfully built a theoretical baseline using Geometric Brownian Motion, dynamically sourcing real-world data to output accurate option pricing via Monte Carlo analysis.

However, the cautionary tale of LTCM proves that relying purely on constant volatility and continuous, smooth paths is dangerous. As we move into the second phase of this project, we intend to abandon the assumption of constant volatility. We will explore advanced stochastic methods, explicitly modifying the underlying math to include discontinuous market shocks (incorporating a Poisson Process for Jump Diffusion), explicitly accounting for the "Crash Risk" that the foundations of modern finance initially failed to see.

\end{document}
